\section{Giới thiệu}

Trong nhiều bài toán học máy hiện đại, mục tiêu của mô hình không chỉ là dự đoán một nhãn rời rạc hoặc một giá trị số, mà còn là sắp xếp một tập các đối tượng theo một thứ tự mong muốn. Bài toán này được gọi là \textit{ranking} hoặc \textit{learning to rank}. Ranking xuất hiện tự nhiên trong nhiều ứng dụng thực tế như xếp hạng kết quả tìm kiếm, gợi ý sản phẩm, sắp xếp tài liệu theo mức độ liên quan, xếp hạng ứng viên, hoặc đánh giá mức độ ưu tiên của các hành động trong một hệ thống ra quyết định.

Khác với bài toán phân loại, trong đó mỗi mẫu thường được gán vào một lớp cụ thể, ranking quan tâm đến quan hệ thứ tự giữa các đối tượng. Ví dụ, trong hệ thống tìm kiếm, điều quan trọng không chỉ là xác định một tài liệu có liên quan hay không, mà còn là tài liệu nào nên được đặt ở vị trí cao hơn trong danh sách kết quả. Do đó, ranking thường được mô hình hoá thông qua các hàm điểm, quan hệ ưu tiên giữa các cặp đối tượng, hoặc các độ đo đánh giá chất lượng của toàn bộ danh sách xếp hạng.

Trong báo cáo này, nhóm tập trung tìm hiểu chủ đề Ranking dựa trên nội dung trong sách \textit{Foundations of Machine Learning}. Mục tiêu chính là trình bày lại các khái niệm nền tảng của ranking bằng tiếng Việt, giải thích ý nghĩa trực giác của các định nghĩa và công thức, đồng thời minh hoạ lý thuyết bằng thực nghiệm lập trình. Báo cáo không chỉ xem ranking như một bài toán ứng dụng, mà còn phân tích các cách tiếp cận lý thuyết phổ biến như \textit{pointwise ranking}, \textit{pairwise ranking} và \textit{listwise ranking}.

Trong ba hướng tiếp cận trên, phần thực nghiệm của báo cáo chọn \textit{pairwise ranking} làm trọng tâm. Lý do là pairwise ranking có mối liên hệ trực tiếp với quan hệ ưu tiên giữa hai đối tượng, ranking loss, và reduction từ ranking sang classification. Bên cạnh đó, hướng tiếp cận này có thể được cài đặt từ đầu bằng các phép toán tuyến tính cơ bản, phù hợp với mục tiêu minh hoạ lý thuyết bằng Python và NumPy. Cụ thể, nhóm xây dựng một mô hình scoring tuyến tính, sinh dữ liệu tổng hợp, tạo các cặp preference, huấn luyện mô hình bằng pairwise hinge loss, và đánh giá kết quả thông qua các độ đo như pairwise accuracy, ranking loss và Kendall tau distance.

Báo cáo được tổ chức như sau. Phần đầu trình bày các khái niệm chuẩn bị cần thiết cho bài toán ranking, bao gồm scoring function, preference relation và ranking loss. Tiếp theo, báo cáo giới thiệu ba hướng tiếp cận chính trong learning to rank: pointwise, pairwise và listwise. Sau đó, phần lý thuyết chính tập trung vào pairwise ranking, cách xây dựng dữ liệu dạng cặp, hàm mất mát pairwise hinge loss và mối liên hệ với bài toán classification. Phần thực nghiệm mô tả quá trình cài đặt mô hình ranking tuyến tính từ đầu, thiết kế dữ liệu tổng hợp, các độ đo đánh giá và phân tích kết quả. Cuối cùng, báo cáo thảo luận một số hướng nghiên cứu hiện đại liên quan đến ranking, đồng thời tổng kết những kết quả đạt được và các hạn chế còn tồn tại.